\chapter{\MakeUppercase{Perancangan Sistem}}

Bab ini menjelaskan rancangan sistem yang diusulkan untuk menyelesaikan permasalahan yang telah dirumuskan pada bab sebelumnya. Perancangan dilakukan berdasarkan landasan teori, hasil studi pustaka, kebutuhan sistem, serta batasan masalah yang telah ditetapkan.

Tujuan utama bab ini adalah menjelaskan bagaimana solusi yang diusulkan akan diwujudkan sebelum dilakukan implementasi. Oleh karena itu, pembahasan pada bab ini harus berfokus pada rancangan sistem, bukan pada hasil implementasi atau hasil pengujian.

Secara umum, bab ini dapat memuat beberapa bagian berikut:

\begin{enumerate}[topsep=0pt,itemsep=0pt,parsep=0pt,leftmargin=*]
\item Gambaran umum sistem yang diusulkan.
\item Analisis kebutuhan sistem, baik kebutuhan fungsional maupun nonfungsional.
\item Arsitektur sistem secara keseluruhan.
\item Diagram blok sistem dan penjelasan fungsi setiap blok.
\item Diagram alir (\textit{flowchart}), \textit{state machine}, \textit{use case}, atau model lain yang sesuai.
\item Perancangan perangkat lunak, algoritma, metode, atau model yang digunakan.
\item Perancangan perangkat keras (jika ada), termasuk hubungan antar komponen dan antarmuka yang digunakan.
\item Perancangan basis data, format data, protokol komunikasi, atau struktur informasi yang digunakan.
\item Perancangan parameter, konfigurasi, dan skenario operasi sistem.
\end{enumerate}

Setiap rancangan yang dibuat harus disertai alasan pemilihannya. Apabila terdapat beberapa alternatif rancangan, jelaskan alasan mengapa rancangan tertentu dipilih untuk digunakan dalam penelitian.

Beberapa hal yang perlu diperhatikan dalam penyusunan Bab Perancangan Sistem adalah sebagai berikut:

\begin{itemize}[topsep=0pt,itemsep=0pt,parsep=0pt,leftmargin=*]
\item Fokus pada \textbf{apa yang akan dibuat dan bagaimana cara kerjanya}, bukan pada hasil implementasi.
\item Gunakan diagram, tabel, atau ilustrasi untuk memperjelas penjelasan sistem.
\item Setiap diagram harus dirujuk dan dijelaskan dalam uraian.
\item Jelaskan alur data, alur proses, masukan, keluaran, serta interaksi antar bagian sistem.
\item Pastikan seluruh rancangan mendukung pencapaian tujuan penelitian dan menjawab rumusan masalah.
\item Hindari pembahasan hasil pengujian, evaluasi, atau analisis performa pada bab ini karena pembahasan tersebut termasuk dalam bab implementasi dan pengujian.
\end{itemize}

Pada akhir bab, pembaca harus dapat memahami struktur sistem yang dirancang, prinsip kerja setiap bagian, serta alasan pemilihan metode, algoritma, maupun konfigurasi yang digunakan sebelum sistem diimplementasikan.
